\section{The replacement thesis and adoption}
\label{sec:adoption}

\subsection{The conditional thesis}

The claim is conditional, and the unconditional version is false. \emph{If} a set of
producers trade with one another regularly, \emph{and} a sufficient share of their inputs
is available from within that set, \emph{then} clearing their obligations against one
another dominates settling each in national currency: the currency each transaction
consumed was only ever a way of moving value the group already collectively possessed.
Where the antecedent fails, so does the conclusion --- a community whose members buy
mostly from outside it cannot clear internally, and no protocol changes that. This is a
mechanism for one identifiable structure, dense recurring mutual trade, and not a general
monetary proposal.

\subsection{The adoption path}

The declared path is continuous trade-credit clearing, entered gradually. Members who
already extend informal trade credit --- pay next month, settle at harvest, run a tab ---
record it here instead of in a notebook. The commercial relationship is unchanged; the
obligation becomes checkable, transferable on the debtor side, insured by the community's
underwriters where their standing reaches (\S\ref{sec:standing}), and eligible for the
consented arbitration of \S\ref{sec:recourse}. A member's first obligation is one they
would have extended anyway, on the terms they would have offered, to somebody they already
trade with: the protocol's contribution at that moment is bookkeeping, and its guarantees
accrue afterwards. Capacity grows as a residue of that ordinary activity --- not applied
for, granted or scored, but what the stakes left behind.

\subsection{Sizing the founding seed}
\label{sec:seedsizing}

A founding community has one number it cannot easily revise upward afterwards: the
external seed. A supply rises only by ceremony (\S\ref{sec:standing}), so the door for a
raise is a second ceremony rather than an edit, and \S\ref{sec:seed-amendments} bounds how
fast it may open. The common error is to size it against annual volume.

The seed is not consumed. It is reserved, released and reserved again, so what it insures
over a year is the seed multiplied by the turnover rate, which settlement terms set and the
seed does not. Measured with one underwriter declaring $10{,}000$ and one debtor: the
debtor's first trade is uninsured, since nobody has staked on them yet, and its settlement
writes the stake; from then on the debtor books the largest amount their residual capacity
insures at term $T$, every epoch, and settles each obligation on its maturity epoch. The
measured year begins one full term after the first such booking, so no term is credited
with the warm-up.

% Auto-generated by scripts/seed-table.py from
% crates/state/examples/seed_table.rs, which measures it against the real
% transition function. Do not hand-edit: change the probe and run
% `just seed-table`. `just seed-table-check` (part of `just ci`) fails the
% build if this file drifts from the measurement it was generated from.
\begin{center}
\begin{tabular}{rrr}
\toprule
settlement term & insured trade per year & $\times$ the seed \\
\midrule
$30$ days & $120{,}000$ & $12.0\times$ \\
$60$ days & $60{,}000$ & $6.0\times$ \\
$90$ days & $40{,}000$ & $4.0\times$ \\
$180$ days & $20{,}000$ & $2.0\times$ \\
$365$ days & $10{,}000$ & $1.0\times$ \\
\bottomrule
\end{tabular}
\end{center}


The shortest row is the floor: a maturity below $\kConstMinMaturityEpochs$ epochs is
refused at acceptance and no governed parameter reaches that field (\S\ref{sec:medium}),
so $365 / \kConstMinMaturityEpochs$ is the ceiling on turnover for any community. Decay
costs the table nothing, by design rather than by the process chosen: an edge never decays
below its live reservation (\S\ref{sec:standing}), so the stake behind an outstanding
obligation cannot fade under it, and the settlement that releases the reservation
re-stakes the edge in the same epoch.

The quantity to found on is therefore \textbf{peak simultaneous insured credit, never
annual volume}: a community expecting $100{,}000$ of trade a year on monthly terms needs a
seed near $10{,}000$, not near $100{,}000$.

The seed also bounds membership. A row is seated by one bond unit of flow held on the
seed's reach and never released (\S\ref{sec:implementation}), so a community seats
$\sum \supply$ over the unit rows and then nobody until a ceremony: at the genesis unit,
one row per $20.00$ of seed, so a founding membership of $M$ needs a seed of at least $M$
units before any credit is sized. The rehearsal (\texttt{the\_seed\_sizing\_rehearsal})
seats founding members by trade, arrives newcomers through members with standing, retires
members at the founder's churn, runs the seed ceremony on the cadence the founder names,
and reports the seats beside the credit: at its default process a seed of $5{,}000$ seats
its $250$ rows and refuses the next thirty-nine epochs into the measured year, while
$20{,}000$ seats every arrival. The rate headroom of \S\ref{sec:seed-amendments} does not
carry over between epochs, so a ceremony a month grows the ceiling by $27\%$ a year and
only a ceremony every epoch compounds at $\beta$.

Who seats a row is a fact about standing, and standing is earned by owing. A seat is held
on its \emph{sponsor's} reach and selling earns none, so the member best placed to bring a
newcomer in --- the shopkeeper letting them run a tab --- can seat nobody until they have
borrowed and repaid themselves, and beyond the one self-act their seat paid for ---
registering guardians, once (\S\ref{sec:implementation}) --- administer their own account
only with a co-signer who has. Onboarding and a member's own administration concentrate
in borrowers and underwriters, and a founder should expect to seat the first members
themselves. A community on a seasonal rhythm sizes the decay ratio before the seed: at the
genesis ratio an unrenewed edge keeps $12\%$ after a quarter and $0.02\%$ after a year
(\S\ref{sec:standing}), and the rehearsal prints the figure beside the seats.

The two errors are not symmetric. Understating the seed costs friction: more obligations
fall to the uninsured tier, which works and is how the network bootstraps. Overstating it
has no later correction, because the one assumption the protocol makes is that
underwriters can bear what they declared (\S\ref{sec:security}). Err low, and re-estimate
against what the ledger measures: utilisation ($\sum$ reserved over $\sum$ supply); the
count of obligations that fell to uninsured for want of capacity --- against a ceiling of
$15{,}000$ with members each wanting $2{,}500$, six members refuse none, ten refuse four,
twenty refuse fourteen and forty refuse thirty-four; and the realised loss rate, the
actuarially correct input and the only one that says whether the declared supply was
priced right.

Unlike the validator set, the underwriter set is \emph{not} frozen at a founding: a
community whose credit could only originate with its founders would wind down as they
aged out, since a stake decays and a supply does not renew itself
(\S\ref{sec:standing}). Members enter the role by an endorsed amendment and leave it by
withdrawing, floored at the flow already committed through them. Opening the role keeps a
community alive without making it richer: for that, an underwriter has to bring something
in from outside, which is a second ceremony (\S\ref{sec:governance}), or the community has
to merge with another.

\paragraph{Who the ceremony is for.} An insured loss is paid in the underwriters' own
obligations (\S\ref{sec:substitution}), so insurance is worth what those obligations buy,
and that is what a seed ceremony decides. A community sizes its seed in underwriters whose
obligations circulate --- a wholesaler its members buy from, a cooperative that sells to
them, a municipality that takes them in fees --- and an underwriter nobody buys from
insures little however much it declares. The first assumption of \S\ref{sec:security} is
the same sentence from the other side.

\paragraph{Why anybody underwrites.} The ledger pays no rent (Property~\ref{prop:no-rent})
and does not compensate the role; four things outside the medium's operation do. The
vote: underwriters are the electorate, and nobody else is (\S\ref{sec:governance}). The
trade: an underwriter is an institution whose members' or customers' credit it wants to
exist --- a wholesaler behind its retailers, a cooperative behind its members --- which is
why suppliers extend terms at their own risk today. The price: a premium is an ordinary
obligation from the insured to the underwriter, which the ledger cannot tell from a trade
and treats as one in every respect, standing included; settled, it stakes the underwriter
on the debtor, bounded by the underwriter's supply like any creditor's stake, so it is no
new lever. The precedent: mutual guarantee societies have filled exactly this role for
decades, on membership fees and public backing.

\subsection{Acceptance, and what actually binds}

The question usually put here is whether a producer will accept an obligation in this
medium for real goods, and the usual answer --- that nothing in the mechanism compels it
--- describes a different mechanism. A sale discharges before it creates
(\S\ref{sec:sale}), so a producer carrying obligations is paid by ceasing to owe, and the
motive is the one that made them a debtor. Measured, netting a debt of $80$ against one's
creditor writes the whole $80$ of stake, where clearing the same $80$ through the cascade
writes none (\S\ref{sec:verification}).

The argument does not reach a member who owes nothing and buys nothing from the community.
For them a claim is what \S\ref{sec:economics} says it is, an expectation of future
delivery, worth taking only if there are things worth buying from members at prices worth
paying. That member is the one the thesis's antecedent already excludes: \textbf{the
binding constraint is the antecedent, not the willingness}. What share of a community's
inputs can be bought from within it is a structural property of that community,
ascertainable before a line of this is deployed, and it decides both whether clearing
dominates and whether anybody has a motive to accept. A community that satisfies it does
not have to be persuaded; one that does not cannot be, by this or any mechanism, and
should not adopt one.

A smaller cost sits beside it: the newcomer's first trade (Remark~\ref{rem:friction}).
Somebody must extend uninsured credit before capacity exists --- a shopkeeper letting a new
neighbour run a tab --- and that is what writes the first stake. The design bounds that
risk and states it; it does not remove it.
